\nivelbasico

\begin{ejercicio}
Una partícula se mueve sobre una circunferencia de radio
\(R=\SI{0.80}{\meter}\) y realiza un desplazamiento angular de
\(135\grado\) en sentido antihorario.

Determine:

\begin{enumerate}[label=\alph*)]
    \item El desplazamiento angular en radianes.
    \item La longitud del arco recorrido.
    \item El signo del desplazamiento angular según la convención usual.
\end{enumerate}
\end{ejercicio}

\begin{respuesta}
\[
\boxed{\Delta\theta=\frac{3\pi}{4}\ \si{\radian}
\approx\SI{2.36}{\radian}}
\]

\[
\boxed{d\approx\SI{1.88}{\meter}}
\]

El desplazamiento angular es positivo.
\end{respuesta}

\begin{explicacion}
Elegimos como positivo el sentido antihorario.

Primero convertimos el ángulo a radianes:

\[
\Delta\theta
=
135\grado
\left(\frac{\pi\ \si{\radian}}{180\grado}\right).
\]

Por tanto:

\[
\boxed{\Delta\theta=\frac{3\pi}{4}\ \si{\radian}}.
\]

Como el radio es constante y no hay inversión del sentido de giro,
la distancia recorrida coincide con la longitud del arco:

\[
d=R|\Delta\theta|.
\]

Sustituyendo:

\[
d
=
(\SI{0.80}{\meter})
\left(\frac{3\pi}{4}\right).
\]

Entonces:

\[
\boxed{d\approx\SI{1.88}{\meter}}.
\]

El signo es positivo porque el desplazamiento angular ocurre en
sentido antihorario.

La distancia recorrida es una magnitud escalar positiva; el signo
pertenece al desplazamiento angular orientado.
\end{explicacion}


\begin{ejercicio}
Una partícula describe una circunferencia de radio
\(R=\SI{0.75}{\meter}\) y completa una revolución en sentido
antihorario.

Determine:

\begin{enumerate}[label=\alph*)]
    \item La distancia recorrida.
    \item El desplazamiento angular acumulado.
    \item El vector desplazamiento entre las posiciones inicial y final.
\end{enumerate}
\end{ejercicio}

\begin{respuesta}
\[
\boxed{d\approx\SI{4.71}{\meter}}
\]

\[
\boxed{\Delta\theta=2\pi\ \si{\radian}}
\]

\[
\boxed{\Delta\vect{r}=\vect{0}}
\]
\end{respuesta}

\begin{explicacion}
La distancia recorrida durante una revolución completa es la
longitud de la circunferencia:

\[
d=2\pi R.
\]

Entonces:

\[
d=2\pi(\SI{0.75}{\meter}).
\]

Por tanto:

\[
\boxed{d\approx\SI{4.71}{\meter}}.
\]

Una revolución antihoraria corresponde a:

\[
\boxed{\Delta\theta=2\pi\ \si{\radian}}.
\]

Para determinar el desplazamiento vectorial utilizamos:

\[
\Delta\vect{r}=\vect{r}_f-\vect{r}_i.
\]

Como la partícula regresa exactamente a su posición inicial:

\[
\vect{r}_f=\vect{r}_i.
\]

Por tanto:

\[
\boxed{\Delta\vect{r}=\vect{0}}.
\]

Este ejercicio muestra que distancia, desplazamiento angular y
desplazamiento vectorial no representan la misma magnitud.
\end{explicacion}


\begin{ejercicio}
Un disco realiza \(80\) revoluciones por minuto con rapidez
angular constante. Un punto de su borde se encuentra a
\(R=\SI{0.25}{\meter}\) del centro.

Determine:

\begin{enumerate}[label=\alph*)]
    \item La frecuencia.
    \item El período.
    \item La rapidez angular.
    \item La rapidez lineal del punto.
\end{enumerate}
\end{ejercicio}

\begin{respuesta}
\[
\boxed{f\approx\SI{1.33}{\hertz}}
\]

\[
\boxed{T=\SI{0.750}{\second}}
\]

\[
\boxed{|\omega|\approx\SI{8.38}{\radian\per\second}}
\]

\[
\boxed{v\approx\SI{2.09}{\meter\per\second}}
\]
\end{respuesta}

\begin{explicacion}
La frecuencia expresa el número de revoluciones por segundo.

Convertimos:

\[
f=\frac{80}{60}\ \si{\per\second}.
\]

Por tanto:

\[
\boxed{f\approx\SI{1.33}{\hertz}}.
\]

El período es:

\[
T=\frac{1}{f}=\frac{60}{80}\ \si{\second}.
\]

Así:

\[
\boxed{T=\SI{0.750}{\second}}.
\]

Como el movimiento es circular uniforme:

\[
|\omega|=2\pi f.
\]

Sustituyendo:

\[
|\omega|=2\pi\left(\frac{80}{60}\right).
\]

Entonces:

\[
\boxed{|\omega|\approx\SI{8.38}{\radian\per\second}}.
\]

La rapidez lineal se relaciona con la rapidez angular mediante:

\[
v=R|\omega|.
\]

Por tanto:

\[
v=(\SI{0.25}{\meter})
(\SI{8.38}{\radian\per\second}).
\]

Así:

\[
\boxed{v\approx\SI{2.09}{\meter\per\second}}.
\]

La rapidez lineal corresponde específicamente al punto situado a
\(\SI{0.25}{\meter}\) del eje.
\end{explicacion}


\begin{ejercicio}
Una rueda gira uniformemente en sentido antihorario y completa
una revolución cada \(\SI{8.00}{\second}\).

Un punto de su borde describe una circunferencia de radio
\(R=\SI{2.00}{\meter}\).

Determine:

\begin{enumerate}[label=\alph*)]
    \item La velocidad angular algebraica.
    \item La rapidez lineal.
    \item La magnitud de la aceleración centrípeta.
    \item La dirección de la aceleración en cualquier punto.
\end{enumerate}
\end{ejercicio}

\begin{respuesta}
\[
\boxed{\omega\approx+\SI{0.785}{\radian\per\second}}
\]

\[
\boxed{v\approx\SI{1.57}{\meter\per\second}}
\]

\[
\boxed{a_c\approx\SI{1.23}{\meter\per\second\squared}}
\]

La aceleración apunta hacia el centro.
\end{respuesta}

\begin{explicacion}
Adoptamos el sentido antihorario como positivo.

En MCU:

\[
\omega=\frac{2\pi}{T}.
\]

Por tanto:

\[
\omega=\frac{2\pi}{8.00}.
\]

Así:

\[
\boxed{\omega\approx+\SI{0.785}{\radian\per\second}}.
\]

La rapidez es:

\[
v=R|\omega|.
\]

Entonces:

\[
v=(2.00)(0.785).
\]

Por tanto:

\[
\boxed{v\approx\SI{1.57}{\meter\per\second}}.
\]

La aceleración centrípeta tiene magnitud:

\[
a_c=R\omega^2.
\]

Sustituyendo:

\[
a_c=(2.00)(0.785)^2.
\]

Así:

\[
\boxed{a_c\approx\SI{1.23}{\meter\per\second\squared}}.
\]

Aunque su magnitud es constante, la dirección de la aceleración
cambia continuamente, pues siempre apunta desde la partícula
hacia el centro de la circunferencia.
\end{explicacion}


\begin{ejercicio}
Un automóvil recorre una curva circular de radio
\(R=\SI{50.0}{\meter}\) con rapidez constante de
\(v=\SI{15.0}{\meter\per\second}\).

Determine:

\begin{enumerate}[label=\alph*)]
    \item La magnitud de su aceleración centrípeta.
    \item Si el vector aceleración es constante.
    \item Qué ocurre con la aceleración centrípeta si se duplica
    la rapidez y se conserva el radio.
\end{enumerate}
\end{ejercicio}

\begin{respuesta}
\[
\boxed{a_c=\SI{4.50}{\meter\per\second\squared}}
\]

El vector aceleración no es constante.

Al duplicar la rapidez:

\[
\boxed{a_c'=4a_c=\SI{18.0}{\meter\per\second\squared}}.
\]
\end{respuesta}

\begin{explicacion}
La aceleración centrípeta tiene magnitud:

\[
a_c=\frac{v^2}{R}.
\]

Entonces:

\[
a_c=\frac{(15.0)^2}{50.0}.
\]

Por tanto:

\[
\boxed{a_c=\SI{4.50}{\meter\per\second\squared}}.
\]

Como el movimiento es circular uniforme, la magnitud de la
aceleración permanece constante.

Sin embargo, el vector aceleración cambia de dirección al
desplazarse el automóvil por la curva.

Por tanto, el vector no es constante.

Si duplicamos la rapidez:

\[
v'=2v,
\]

la nueva aceleración es:

\[
a_c'=\frac{(2v)^2}{R}.
\]

Desarrollando:

\[
a_c'=4\frac{v^2}{R}=4a_c.
\]

Entonces:

\[
a_c'=4(4.50).
\]

Así:

\[
\boxed{a_c'=\SI{18.0}{\meter\per\second\squared}}.
\]

El resultado demuestra que la aceleración centrípeta depende
cuadráticamente de la rapidez.
\end{explicacion}


\begin{ejercicio}
Una partícula describe un movimiento circular uniforme en sentido
antihorario.

En cierto instante se encuentra en el punto más a la derecha
de su trayectoria.

El centro de la circunferencia coincide con el origen y el
sistema cartesiano tiene \(+x\) hacia la derecha y \(+y\)
hacia arriba.

Determine la dirección y el sentido de:

\begin{enumerate}[label=\alph*)]
    \item El vector posición.
    \item El vector velocidad.
    \item El vector aceleración.
\end{enumerate}

Explique por qué la velocidad y la aceleración son perpendiculares.
\end{ejercicio}

\begin{respuesta}
\[
\vect{r}=R\ihat,
\]

\[
\vect{v}=v\jhat,
\]

\[
\vect{a}=-a_c\ihat.
\]

La velocidad apunta hacia \(+y\) y la aceleración hacia \(-x\).
\end{respuesta}

\begin{explicacion}
En el punto más a la derecha de la circunferencia:

\[
x=R,
\qquad
y=0.
\]

Por tanto:

\[
\boxed{\vect{r}=R\ihat}.
\]

La velocidad instantánea es tangente a la trayectoria.

Como el giro es antihorario, en ese punto la velocidad
apunta verticalmente hacia arriba:

\[
\boxed{\vect{v}=v\jhat}.
\]

En MCU, la aceleración es centrípeta y apunta hacia el centro.

Desde el punto derecho de la circunferencia, el centro
se encuentra hacia la izquierda.

Por tanto:

\[
\boxed{\vect{a}=-a_c\ihat}.
\]

Como:

\[
\ihat\pdot\jhat=0,
\]

obtenemos:

\[
\vect{v}\pdot\vect{a}=0.
\]

Esto significa que la aceleración modifica la dirección del
vector velocidad, pero no su magnitud.
\end{explicacion}


\nivelintermedio

\begin{ejercicio}
Una partícula describe una circunferencia de radio
\(R=\SI{1.50}{\meter}\) alrededor del origen.

En un instante su posición angular es \(120\grado\),
medida desde el eje \(+x\), y su velocidad angular es:

\[
\omega=-\SI{2.00}{\radian\per\second}.
\]

Determine:

\begin{enumerate}[label=\alph*)]
    \item El vector posición.
    \item El vector velocidad.
    \item El vector aceleración.
    \item La rapidez y la magnitud de la aceleración.
\end{enumerate}

Considere MCU.
\end{ejercicio}

\begin{respuesta}
\[
\boxed{
\vect{r}
\approx(-0.750\ihat+1.30\jhat)\,\si{\meter}
}
\]

\[
\boxed{
\vect{v}
\approx(2.60\ihat+1.50\jhat)\,
\si{\meter\per\second}
}
\]

\[
\boxed{
\vect{a}
\approx(3.00\ihat-5.20\jhat)\,
\si{\meter\per\second\squared}
}
\]

\[
\boxed{v=\SI{3.00}{\meter\per\second}},
\qquad
\boxed{a=\SI{6.00}{\meter\per\second\squared}}.
\]
\end{respuesta}

\begin{explicacion}
Tomamos el centro como origen, \(+x\) hacia la derecha y
\(+y\) hacia arriba. El signo negativo de \(\omega\) indica
giro horario.

La posición cartesiana es:

\[
x=R\cos\theta,
\qquad
y=R\sen\theta.
\]

Con \(\theta=120\grado\):

\[
x=(1.50)\cos120\grado=-0.750,
\]

\[
y=(1.50)\sen120\grado\approx1.30.
\]

Por tanto:

\[
\boxed{
\vect{r}\approx(-0.750\ihat+1.30\jhat)\,\si{\meter}
}.
\]

Para la velocidad utilizamos:

\[
v_x=-R\omega\sen\theta,
\qquad
v_y=R\omega\cos\theta.
\]

Entonces:

\[
v_x=-(1.50)(-2.00)\sen120\grado
\approx+\SI{2.60}{\meter\per\second},
\]

\[
v_y=(1.50)(-2.00)\cos120\grado
=+\SI{1.50}{\meter\per\second}.
\]

Así:

\[
\boxed{
\vect{v}\approx(2.60\ihat+1.50\jhat)\,
\si{\meter\per\second}
}.
\]

Como el movimiento es circular uniforme:

\[
\vect{a}=-\omega^2\vect{r}.
\]

Trabajando por componentes:

\[
a_x=-\omega^2x,
\qquad
a_y=-\omega^2y.
\]

Por tanto:

\[
a_x=-(2.00)^2(-0.750)
=+\SI{3.00}{\meter\per\second\squared},
\]

\[
a_y=-(2.00)^2(1.50\sen120\grado)
\approx-\SI{5.20}{\meter\per\second\squared}.
\]

Entonces:

\[
\boxed{
\vect{a}\approx(3.00\ihat-5.20\jhat)\,
\si{\meter\per\second\squared}
}.
\]

Las magnitudes también pueden obtenerse mediante:

\[
v=R|\omega|=(1.50)(2.00)
=\SI{3.00}{\meter\per\second},
\]

\[
a=R\omega^2=(1.50)(2.00)^2
=\SI{6.00}{\meter\per\second\squared}.
\]

La aceleración apunta hacia el centro, mientras que la velocidad
es tangente a la circunferencia y corresponde al giro horario.
\end{explicacion}


\begin{ejercicio}
Una rueda inicialmente en reposo comienza a girar en sentido
antihorario con aceleración angular constante:

\[
\alpha=\SI{1.50}{\radian\per\second\squared}.
\]

Un punto de su borde está a \(R=\SI{0.400}{\meter}\)
del eje de rotación.

Después de \(\SI{4.00}{\second}\), determine:

\begin{enumerate}[label=\alph*)]
    \item La velocidad angular.
    \item El desplazamiento angular acumulado.
    \item El número de revoluciones.
    \item La rapidez lineal.
    \item Las magnitudes de las aceleraciones tangencial y centrípeta.
\end{enumerate}
\end{ejercicio}

\begin{respuesta}
\[
\boxed{\omega=\SI{6.00}{\radian\per\second}}
\]

\[
\boxed{\Delta\theta=\SI{12.0}{\radian}}
\]

\[
\boxed{N\approx1.91\text{ revoluciones}}
\]

\[
\boxed{v=\SI{2.40}{\meter\per\second}}
\]

\[
\boxed{a_t=\SI{0.600}{\meter\per\second\squared}}
\]

\[
\boxed{a_c=\SI{14.4}{\meter\per\second\squared}}
\]
\end{respuesta}

\begin{explicacion}
Adoptamos el giro antihorario como positivo. Dado que la
rueda parte del reposo:

\[
\omega_0=0.
\]

Como la aceleración angular es constante:

\[
\omega=\omega_0+\alpha t.
\]

Entonces:

\[
\omega=(1.50)(4.00)
=\boxed{\SI{6.00}{\radian\per\second}}.
\]

El desplazamiento angular es:

\[
\Delta\theta=\omega_0t+\frac12\alpha t^2.
\]

Sustituyendo:

\[
\Delta\theta=\frac12(1.50)(4.00)^2.
\]

Por tanto:

\[
\boxed{\Delta\theta=\SI{12.0}{\radian}}.
\]

Para convertir a revoluciones:

\[
N=\frac{\Delta\theta}{2\pi}.
\]

Así:

\[
\boxed{N\approx1.91\text{ revoluciones}}.
\]

La rapidez del punto situado en el borde es:

\[
v=R|\omega|.
\]

Entonces:

\[
v=(0.400)(6.00)
=\boxed{\SI{2.40}{\meter\per\second}}.
\]

La aceleración tangencial es:

\[
a_t=R\alpha=(0.400)(1.50).
\]

Por tanto:

\[
\boxed{a_t=\SI{0.600}{\meter\per\second\squared}}.
\]

La aceleración centrípeta:

\[
a_c=R\omega^2=(0.400)(6.00)^2.
\]

Así:

\[
\boxed{a_c=\SI{14.4}{\meter\per\second\squared}}.
\]

La aceleración tangencial hace aumentar la rapidez del punto.
La aceleración centrípeta modifica continuamente la dirección
de su velocidad.
\end{explicacion}


\begin{ejercicio}
Una rueda gira inicialmente en sentido antihorario con:

\[
\omega_0=\SI{10.0}{\radian\per\second}.
\]

A partir de cierto instante experimenta una aceleración
angular constante:

\[
\alpha=-\SI{2.50}{\radian\per\second\squared}.
\]

Determine:

\begin{enumerate}[label=\alph*)]
    \item El tiempo necesario para detenerse momentáneamente.
    \item El desplazamiento angular hasta ese instante.
    \item El número de revoluciones realizadas durante el frenado.
    \item Su velocidad angular a los \(\SI{6.00}{\second}\),
    si se mantiene la misma aceleración.
    \item Su sentido de giro en ese último instante.
\end{enumerate}
\end{ejercicio}

\begin{respuesta}
\[
\boxed{t_{\mathrm{det}}=\SI{4.00}{\second}}
\]

\[
\boxed{\Delta\theta=\SI{20.0}{\radian}}
\]

\[
\boxed{N\approx3.18\text{ revoluciones}}
\]

\[
\boxed{\omega(6.00)=-\SI{5.00}{\radian\per\second}}
\]

A los seis segundos gira en sentido horario.
\end{respuesta}

\begin{explicacion}
Tomamos el giro antihorario como positivo.

Como la aceleración angular es constante:

\[
\omega=\omega_0+\alpha t.
\]

En el instante de detención:

\[
\omega=0.
\]

Por tanto:

\[
0=10.0-2.50t.
\]

Despejando:

\[
\boxed{t_{\mathrm{det}}=\SI{4.00}{\second}}.
\]

El desplazamiento angular hasta detenerse es:

\[
\Delta\theta=\omega_0t+\frac12\alpha t^2.
\]

Sustituyendo:

\[
\Delta\theta
=10.0(4.00)+\frac12(-2.50)(4.00)^2.
\]

Entonces:

\[
\Delta\theta=40.0-20.0.
\]

Así:

\[
\boxed{\Delta\theta=\SI{20.0}{\radian}}.
\]

El número de revoluciones realizadas es:

\[
N=\frac{20.0}{2\pi}.
\]

Por tanto:

\[
\boxed{N\approx3.18\text{ revoluciones}}.
\]

Si la aceleración sigue actuando después de la detención,
el movimiento cambia de sentido.

Para \(t=\SI{6.00}{\second}\):

\[
\omega=10.0-2.50(6.00).
\]

Por tanto:

\[
\boxed{\omega=-\SI{5.00}{\radian\per\second}}.
\]

El signo negativo indica giro horario. Durante los primeros
cuatro segundos disminuye la rapidez de giro; después aumenta,
pero en el sentido contrario.

En el instante de inversión, la velocidad angular es cero,
aunque la aceleración angular sigue siendo distinta de cero.
\end{explicacion}


\begin{ejercicio}
Dos puntos \(A\) y \(B\) pertenecen a un mismo disco rígido
que gira uniformemente alrededor de un eje fijo.

Sus distancias al eje son:

\[
R_A=\SI{0.100}{\meter},
\qquad
R_B=\SI{0.300}{\meter}.
\]

La rapidez angular del disco es:

\[
|\omega|=\SI{12.0}{\radian\per\second}.
\]

Determine:

\begin{enumerate}[label=\alph*)]
    \item La rapidez lineal de cada punto.
    \item La aceleración centrípeta de cada punto.
    \item Las razones \(v_B/v_A\) y \(a_{cB}/a_{cA}\).
    \item Si sus períodos son iguales o diferentes.
\end{enumerate}
\end{ejercicio}

\begin{respuesta}
\[
v_A=\SI{1.20}{\meter\per\second},
\qquad
v_B=\SI{3.60}{\meter\per\second},
\]

\[
a_{cA}=\SI{14.4}{\meter\per\second\squared},
\qquad
a_{cB}=\SI{43.2}{\meter\per\second\squared},
\]

\[
\boxed{\frac{v_B}{v_A}=3},
\qquad
\boxed{\frac{a_{cB}}{a_{cA}}=3}.
\]

Ambos puntos tienen el mismo período.
\end{respuesta}

\begin{explicacion}
Como los puntos pertenecen al mismo disco rígido y éste gira
alrededor de un eje fijo, comparten la velocidad angular.

La rapidez lineal es:

\[
v=R|\omega|.
\]

Para el punto \(A\):

\[
v_A=(0.100)(12.0)
=\boxed{\SI{1.20}{\meter\per\second}}.
\]

Para el punto \(B\):

\[
v_B=(0.300)(12.0)
=\boxed{\SI{3.60}{\meter\per\second}}.
\]

La aceleración centrípeta es:

\[
a_c=R\omega^2.
\]

Para \(A\):

\[
a_{cA}=(0.100)(12.0)^2
=\boxed{\SI{14.4}{\meter\per\second\squared}}.
\]

Para \(B\):

\[
a_{cB}=(0.300)(12.0)^2
=\boxed{\SI{43.2}{\meter\per\second\squared}}.
\]

Como \(|\omega|\) es común:

\[
\frac{v_B}{v_A}=\frac{R_B}{R_A}=3,
\]

\[
\frac{a_{cB}}{a_{cA}}=\frac{R_B}{R_A}=3.
\]

El período depende de la rapidez angular compartida:

\[
T=\frac{2\pi}{|\omega|}.
\]

Por tanto, ambos puntos completan cada revolución en el
mismo tiempo.

El resultado muestra que compartir velocidad angular no
significa compartir rapidez lineal ni aceleración centrípeta.
\end{explicacion}


\begin{ejercicio}
Una partícula se desplaza sobre una circunferencia de radio
\(R=\SI{2.00}{\meter}\).

En un instante determinado:

\[
\theta=30.0\grado,
\qquad
\omega=\SI{3.00}{\radian\per\second},
\]

\[
\alpha=-\SI{1.00}{\radian\per\second\squared}.
\]

El centro coincide con el origen y los ángulos positivos se
miden en sentido antihorario.

Determine:

\begin{enumerate}[label=\alph*)]
    \item Las componentes rectangulares de la velocidad.
    \item Las componentes rectangulares de la aceleración tangencial.
    \item Las componentes rectangulares de la aceleración centrípeta.
    \item El vector aceleración total y su magnitud.
    \item Si la rapidez aumenta o disminuye.
\end{enumerate}
\end{ejercicio}

\begin{respuesta}
\[
\boxed{
\vect{v}
=(-3.00\ihat+5.20\jhat)\,
\si{\meter\per\second}
}
\]

\[
\boxed{
\vect{a}_t
=(1.00\ihat-1.73\jhat)\,
\si{\meter\per\second\squared}
}
\]

\[
\boxed{
\vect{a}_c
=(-15.6\ihat-9.00\jhat)\,
\si{\meter\per\second\squared}
}
\]

\[
\boxed{
\vect{a}
\approx(-14.6\ihat-10.7\jhat)\,
\si{\meter\per\second\squared}
}
\]

\[
\boxed{a\approx\SI{18.1}{\meter\per\second\squared}}
\]

La rapidez disminuye.
\end{respuesta}

\begin{explicacion}
Utilizamos ejes cartesianos fijos con el centro de la
circunferencia en el origen.

Las componentes de velocidad son:

\[
v_x=-R\omega\sen\theta,
\qquad
v_y=R\omega\cos\theta.
\]

Sustituyendo:

\[
v_x=-(2.00)(3.00)\sen30\grado=-3.00,
\]

\[
v_y=(2.00)(3.00)\cos30\grado\approx5.20.
\]

Por tanto:

\[
\boxed{
\vect{v}
\approx(-3.00\ihat+5.20\jhat)\,
\si{\meter\per\second}
}.
\]

Las componentes de aceleración tangencial son:

\[
a_{tx}=-R\alpha\sen\theta,
\qquad
a_{ty}=R\alpha\cos\theta.
\]

Entonces:

\[
a_{tx}=-(2.00)(-1.00)\sen30\grado=+1.00,
\]

\[
a_{ty}=(2.00)(-1.00)\cos30\grado\approx-1.73.
\]

Así:

\[
\boxed{
\vect{a}_t
\approx(1.00\ihat-1.73\jhat)\,
\si{\meter\per\second\squared}
}.
\]

La aceleración centrípeta se obtiene mediante:

\[
a_{cx}=-R\omega^2\cos\theta,
\]

\[
a_{cy}=-R\omega^2\sen\theta.
\]

Por tanto:

\[
a_{cx}=-(2.00)(3.00)^2\cos30\grado\approx-15.6,
\]

\[
a_{cy}=-(2.00)(3.00)^2\sen30\grado=-9.00.
\]

Entonces:

\[
\boxed{
\vect{a}_c
\approx(-15.6\ihat-9.00\jhat)\,
\si{\meter\per\second\squared}
}.
\]

Sumamos las componentes:

\[
a_x=a_{tx}+a_{cx},
\qquad
a_y=a_{ty}+a_{cy}.
\]

Así:

\[
a_x\approx-14.6,
\qquad
a_y\approx-10.7.
\]

Por tanto:

\[
\boxed{
\vect{a}
\approx(-14.6\ihat-10.7\jhat)\,
\si{\meter\per\second\squared}
}.
\]

La magnitud también puede calcularse a partir de las dos
componentes perpendiculares:

\[
a_t=R|\alpha|=2.00,
\qquad
a_c=R\omega^2=18.0.
\]

Entonces:

\[
a=\sqrt{a_t^2+a_c^2}
=\sqrt{(2.00)^2+(18.0)^2}.
\]

Así:

\[
\boxed{a\approx\SI{18.1}{\meter\per\second\squared}}.
\]

Como \(\omega>0\) y \(\alpha<0\), la velocidad angular
algebraica todavía corresponde al giro antihorario, pero
su magnitud está disminuyendo.

En consecuencia, la aceleración tangencial apunta en
sentido contrario a la velocidad.
\end{explicacion}


\begin{ejercicio}
El registro de un sistema giratorio muestra que su velocidad
angular, expresada en \(\si{\radian\per\second}\), cambia con el
tiempo según la siguiente gráfica.

\begin{center}
\begin{tikzpicture}
\begin{axis}[
    width=0.78\textwidth,
    height=5.8cm,
    axis lines=left,
    xlabel={\(t\) (\si{\second})},
    ylabel={\(\omega\) (\si{\radian\per\second})},
    xmin=0,
    xmax=7,
    ymin=-1,
    ymax=8,
    xtick={0,2,4,6},
    ytick={0,2,4,6,8},
    grid=major
]
\addplot[
    AzulPrincipal,
    very thick,
    mark=*,
]
coordinates {(0,2) (4,6) (6,6)};
\end{axis}
\end{tikzpicture}
\end{center}

Determine:

\begin{enumerate}[label=\alph*)]
    \item La aceleración angular entre \(0\) y \(\SI{4}{\second}\).
    \item La aceleración angular entre \(\SI{4}{\second}\) y
    \(\SI{6}{\second}\).
    \item El desplazamiento angular acumulado entre \(0\) y
    \(\SI{6}{\second}\).
    \item El número de revoluciones correspondiente.
\end{enumerate}
\end{ejercicio}

\begin{respuesta}
\[
\boxed{\alpha_{0-4}=\SI{1.00}{\radian\per\second\squared}}
\]

\[
\boxed{\alpha_{4-6}=0}
\]

\[
\boxed{\Delta\theta=\SI{28.0}{\radian}}
\]

\[
\boxed{N\approx4.46\text{ revoluciones}}
\]
\end{respuesta}

\begin{explicacion}
La pendiente de una gráfica \(\omega-t\) representa la
aceleración angular.

Durante el primer intervalo:

\[
\alpha_{0-4}=\frac{\Delta\omega}{\Delta t}.
\]

Entonces:

\[
\alpha_{0-4}=\frac{6-2}{4-0}.
\]

Por tanto:

\[
\boxed{\alpha_{0-4}=\SI{1.00}{\radian\per\second\squared}}.
\]

Entre \(t=\SI{4}{\second}\) y \(t=\SI{6}{\second}\),
la gráfica es horizontal.

Por tanto:

\[
\boxed{\alpha_{4-6}=0}.
\]

El desplazamiento angular corresponde al área algebraica
bajo la gráfica \(\omega-t\).

Entre \(0\) y \(\SI{4}{\second}\), el área del trapecio es:

\[
\Delta\theta_1
=\frac{\omega_i+\omega_f}{2}\Delta t.
\]

Sustituyendo:

\[
\Delta\theta_1=\frac{2+6}{2}(4)
=\SI{16.0}{\radian}.
\]

Entre \(\SI{4}{\second}\) y \(\SI{6}{\second}\),
el área es rectangular:

\[
\Delta\theta_2=\omega\Delta t=6(2).
\]

Entonces:

\[
\Delta\theta_2=\SI{12.0}{\radian}.
\]

El desplazamiento angular total es:

\[
\Delta\theta=\Delta\theta_1+\Delta\theta_2.
\]

Por tanto:

\[
\boxed{\Delta\theta=\SI{28.0}{\radian}}.
\]

Como todo el giro ocurre en el mismo sentido, el número
de revoluciones es:

\[
N=\frac{28.0}{2\pi}.
\]

Así:

\[
\boxed{N\approx4.46\text{ revoluciones}}.
\]

El primer tramo corresponde a un aumento uniforme de
velocidad angular; el segundo, a movimiento circular uniforme.
\end{explicacion}


\nivelavanzado

\begin{ejercicio}
Una partícula se desplaza sobre una circunferencia fija de
radio \(R\), centrada en el origen del plano \(xy\).

Su posición angular está dada por:

\[
\theta(t)=\theta_0+\omega_0t+\frac12\alpha t^2,
\]

donde \(\alpha\) es constante.

Sin utilizar directamente fórmulas conocidas para las
componentes de aceleración:

\begin{enumerate}[label=\alph*)]
    \item Deduzca las componentes rectangulares de la posición.
    \item Obtenga las componentes de velocidad.
    \item Obtenga las componentes de aceleración.
    \item Identifique la contribución tangencial y la centrípeta.
    \item Demuestre que son perpendiculares.
    \item Compruebe que el resultado recupera el MCU cuando
    \(\alpha=0\).
\end{enumerate}
\end{ejercicio}

\begin{respuesta}
\[
\boxed{
\vect{r}=R\cos\theta\,\ihat+R\sen\theta\,\jhat
}
\]

\[
\boxed{
\vect{v}=-R\omega\sen\theta\,\ihat
+R\omega\cos\theta\,\jhat
}
\]

\[
\boxed{
\vect{a}
=
(-R\alpha\sen\theta-R\omega^2\cos\theta)\ihat
+
(R\alpha\cos\theta-R\omega^2\sen\theta)\jhat
}
\]

\[
\boxed{
\vect{a}_t\pdot\vect{a}_c=0
}
\]
\end{respuesta}

\begin{explicacion}
La partícula permanece sobre una circunferencia de radio
constante, de modo que:

\[
x=R\cos\theta,
\qquad
y=R\sen\theta.
\]

Entonces:

\[
\boxed{
\vect{r}
=R\cos\theta\,\ihat+R\sen\theta\,\jhat
}.
\]

La posición angular depende del tiempo. Por tanto, al derivar
debemos aplicar la regla de la cadena.

Primero:

\[
\omega=\deriv{\theta}{t}=\omega_0+\alpha t.
\]

Para la velocidad:

\[
v_x=\deriv{x}{t}
=-R\sen\theta\,\deriv{\theta}{t}.
\]

Entonces:

\[
v_x=-R\omega\sen\theta.
\]

De manera semejante:

\[
v_y=\deriv{y}{t}
=R\omega\cos\theta.
\]

Por tanto:

\[
\boxed{
\vect{v}
=-R\omega\sen\theta\,\ihat
+R\omega\cos\theta\,\jhat
}.
\]

Como \(\omega\) también cambia con el tiempo, para obtener
la aceleración debemos aplicar la regla del producto.

En el eje \(x\):

\[
a_x=\deriv{}{t}(-R\omega\sen\theta).
\]

Entonces:

\[
a_x=-R\alpha\sen\theta-R\omega^2\cos\theta.
\]

En el eje \(y\):

\[
a_y=\deriv{}{t}(R\omega\cos\theta).
\]

Por tanto:

\[
a_y=R\alpha\cos\theta-R\omega^2\sen\theta.
\]

Así:

\[
\boxed{
\vect{a}
=
(-R\alpha\sen\theta-R\omega^2\cos\theta)\ihat
+
(R\alpha\cos\theta-R\omega^2\sen\theta)\jhat
}.
\]

Agrupamos los términos que contienen \(\alpha\):

\[
\vect{a}_t
=
R\alpha(-\sen\theta\,\ihat+\cos\theta\,\jhat).
\]

Su dirección es tangencial.

Los términos que contienen \(\omega^2\) son:

\[
\vect{a}_c
=
-R\omega^2(\cos\theta\,\ihat+\sen\theta\,\jhat).
\]

Estos apuntan radialmente hacia el centro.

Para comprobar perpendicularidad calculamos el producto
escalar:

\[
\vect{a}_t\pdot\vect{a}_c
=
R^2\alpha\omega^2
\left(
\sen\theta\cos\theta
-
\cos\theta\sen\theta
\right).
\]

Por tanto:

\[
\boxed{\vect{a}_t\pdot\vect{a}_c=0}.
\]

Finalmente, si:

\[
\alpha=0,
\]

desaparece la aceleración tangencial:

\[
\vect{a}_t=\vect{0}.
\]

La aceleración queda:

\[
\vect{a}
=-R\omega^2(\cos\theta\,\ihat+\sen\theta\,\jhat).
\]

Como:

\[
\vect{r}=R(\cos\theta\,\ihat+\sen\theta\,\jhat),
\]

resulta:

\[
\boxed{\vect{a}=-\omega^2\vect{r}}.
\]

Es exactamente la aceleración centrípeta del MCU.
\end{explicacion}


\begin{ejercicio}
Una partícula describe una circunferencia de radio
\(R=\SI{1.20}{\meter}\) en el plano \(xy\), con centro
en el origen.

Su posición angular se modela mediante:

\[
\theta(t)
=
\frac{\pi}{6}+2t-\frac12t^2,
\]

donde los términos numéricos se entienden expresados en las
unidades SI correspondientes.

Determine:

\begin{enumerate}[label=\alph*)]
    \item El instante en que la partícula cambia de sentido.
    \item Su posición angular acumulada en ese instante.
    \item Su vector posición en el instante de cambio.
    \item Su vector velocidad en ese instante.
    \item Su vector aceleración en ese instante.
    \item Si la aceleración centrípeta desaparece cuando cambia
    el sentido.
\end{enumerate}
\end{ejercicio}

\begin{respuesta}
\[
\boxed{t_c=\SI{2.00}{\second}}
\]

\[
\boxed{\theta(t_c)=\frac{\pi}{6}+2\ \si{\radian}}
\]

\[
\boxed{
\vect{r}(t_c)
\approx(-0.979\ihat+0.693\jhat)\,\si{\meter}
}
\]

\[
\boxed{\vect{v}(t_c)=\vect{0}}
\]

\[
\boxed{
\vect{a}(t_c)
\approx(0.693\ihat+0.979\jhat)\,
\si{\meter\per\second\squared}
}
\]

La aceleración centrípeta es cero en ese instante.
\end{respuesta}

\begin{explicacion}
La velocidad angular es la derivada de la posición angular:

\[
\omega(t)=\deriv{\theta}{t}=2-t.
\]

La aceleración angular es:

\[
\alpha=\deriv{\omega}{t}
=-\SI{1.00}{\radian\per\second\squared}.
\]

El cambio de sentido ocurre cuando la velocidad angular
pasa por cero:

\[
0=2-t_c.
\]

Por tanto:

\[
\boxed{t_c=\SI{2.00}{\second}}.
\]

La posición angular acumulada es:

\[
\theta(2)
=\frac{\pi}{6}+2(2)-\frac12(2)^2.
\]

Así:

\[
\boxed{\theta(2)=\frac{\pi}{6}+2\ \si{\radian}}.
\]

Su valor aproximado es:

\[
\theta(2)\approx\SI{2.524}{\radian}.
\]

La posición cartesiana se obtiene mediante:

\[
x=R\cos\theta,
\qquad
y=R\sen\theta.
\]

Entonces:

\[
x=(1.20)\cos(2.524)\approx-0.979,
\]

\[
y=(1.20)\sen(2.524)\approx0.693.
\]

Por tanto:

\[
\boxed{
\vect{r}(2)
\approx(-0.979\ihat+0.693\jhat)\,\si{\meter}
}.
\]

Como:

\[
\omega(2)=0,
\]

la rapidez lineal es:

\[
v=R|\omega|=0.
\]

En consecuencia:

\[
\boxed{\vect{v}(2)=\vect{0}}.
\]

La aceleración centrípeta también desaparece:

\[
a_c=R\omega^2=0.
\]

Sin embargo, la aceleración angular continúa siendo:

\[
\alpha=-\SI{1.00}{\radian\per\second\squared}.
\]

La aceleración tangencial es:

\[
\vect{a}_t
=
R\alpha
(-\sen\theta\,\ihat+\cos\theta\,\jhat).
\]

Sustituyendo:

\[
\vect{a}_t
=
(1.20)(-1.00)
[-\sen(2.524)\,\ihat+\cos(2.524)\,\jhat].
\]

Por tanto:

\[
\boxed{
\vect{a}(2)
\approx(0.693\ihat+0.979\jhat)\,
\si{\meter\per\second\squared}
}.
\]

En ese instante la aceleración es exclusivamente tangencial.

La partícula está momentáneamente en reposo, pero no tiene
aceleración nula. La aceleración tangencial hace que comience
a girar en el sentido contrario.
\end{explicacion}


\begin{ejercicio}
Un disco rígido gira alrededor de un eje fijo. Su rapidez
angular cambia con el tiempo y se desea determinar si su
aceleración angular es constante.

Durante un experimento se registran los siguientes datos:

\begin{center}
\begin{tabular}{c|ccccc}
\hline
\(t\) (\si{\second}) & 0 & 1 & 2 & 3 & 4 \\
\hline
\(\omega\) (\si{\radian\per\second}) & 1 & 3 & 5 & 7 & 9 \\
\hline
\end{tabular}
\end{center}

Un punto del disco se encuentra a
\(R=\SI{0.250}{\meter}\) del eje.

Determine:

\begin{enumerate}[label=\alph*)]
    \item Si los datos son compatibles con aceleración angular constante.
    \item La función \(\omega(t)\) correspondiente.
    \item El desplazamiento angular entre \(0\) y
    \(\SI{4.00}{\second}\).
    \item Las magnitudes de la aceleración tangencial y centrípeta
    en \(t=\SI{4.00}{\second}\).
    \item La magnitud de la aceleración total en ese instante.
\end{enumerate}
\end{ejercicio}

\begin{respuesta}
Los datos son compatibles con aceleración angular constante:

\[
\boxed{\alpha=\SI{2.00}{\radian\per\second\squared}}
\]

\[
\boxed{\omega(t)=1+2t}
\]

\[
\boxed{\Delta\theta=\SI{20.0}{\radian}}
\]

\[
\boxed{a_t=\SI{0.500}{\meter\per\second\squared}}
\]

\[
\boxed{a_c=\SI{20.25}{\meter\per\second\squared}}
\]

\[
\boxed{a\approx\SI{20.3}{\meter\per\second\squared}}
\]
\end{respuesta}

\begin{explicacion}
La aceleración angular media entre dos mediciones consecutivas es:

\[
\bar{\alpha}=\frac{\Delta\omega}{\Delta t}.
\]

Entre \(t=0\) y \(t=\SI{1}{\second}\):

\[
\bar{\alpha}=\frac{3-1}{1-0}=2.
\]

Entre \(t=\SI{1}{\second}\) y \(t=\SI{2}{\second}\):

\[
\bar{\alpha}=\frac{5-3}{2-1}=2.
\]

Los demás intervalos proporcionan el mismo resultado.

Por tanto, los datos son compatibles con:

\[
\boxed{\alpha=\SI{2.00}{\radian\per\second\squared}}.
\]

Como:

\[
\omega_0=\SI{1.00}{\radian\per\second},
\]

la ecuación angular es:

\[
\omega=\omega_0+\alpha t.
\]

Así:

\[
\boxed{\omega(t)=1+2t},
\]

con \(\omega\) en \(\si{\radian\per\second}\)
y \(t\) en segundos.

El desplazamiento angular es:

\[
\Delta\theta=\omega_0t+\frac12\alpha t^2.
\]

Para \(t=\SI{4.00}{\second}\):

\[
\Delta\theta=1(4)+\frac12(2)(4)^2.
\]

Entonces:

\[
\boxed{\Delta\theta=\SI{20.0}{\radian}}.
\]

La aceleración tangencial es:

\[
a_t=R\alpha.
\]

Por tanto:

\[
a_t=(0.250)(2.00)
=\boxed{\SI{0.500}{\meter\per\second\squared}}.
\]

En \(t=\SI{4.00}{\second}\):

\[
\omega=1+2(4)=\SI{9.00}{\radian\per\second}.
\]

La aceleración centrípeta es:

\[
a_c=R\omega^2.
\]

Entonces:

\[
a_c=(0.250)(9.00)^2
=\boxed{\SI{20.25}{\meter\per\second\squared}}.
\]

Como las aceleraciones tangencial y centrípeta son
perpendiculares:

\[
a=\sqrt{a_t^2+a_c^2}.
\]

Por tanto:

\[
a=\sqrt{(0.500)^2+(20.25)^2}.
\]

Así:

\[
\boxed{a\approx\SI{20.3}{\meter\per\second\squared}}.
\]

El resultado evidencia que una aceleración angular constante
no implica aceleración lineal total constante.

A medida que aumenta \(\omega\), aumenta la componente
centrípeta.
\end{explicacion}


\begin{ejercicio}
Un estudiante afirma:

\begin{quote}
«Si una partícula describe una circunferencia y su aceleración
tiene magnitud constante, necesariamente realiza movimiento
circular uniforme».
\end{quote}

Analice si la afirmación es correcta.

Para justificar su respuesta, considere una partícula de radio
constante \(R\) y determine si puede tener rapidez variable
mientras la magnitud de su aceleración total permanece constante.

Si es posible, construya matemáticamente un ejemplo mediante
una relación entre \(\omega\) y \(\alpha\).
\end{ejercicio}

\begin{respuesta}
La afirmación es falsa.

Es posible tener:

\[
\boxed{a=\text{constante}}
\]

con:

\[
\boxed{\omega\neq\text{constante}}.
\]

Una condición posible es:

\[
\boxed{
\alpha^2+\omega^4=C
}
\]

para una constante positiva \(C\), con las unidades
correspondientes.
\end{respuesta}

\begin{explicacion}
En una trayectoria circular de radio constante, la magnitud
de la aceleración total es:

\[
a=\sqrt{a_t^2+a_c^2}.
\]

Como:

\[
a_t=R\alpha,
\]

y:

\[
a_c=R\omega^2,
\]

tenemos:

\[
a=R\sqrt{\alpha^2+\omega^4}.
\]

Si deseamos que la magnitud total sea una constante positiva
\(a_0\), debe cumplirse:

\[
a_0=R\sqrt{\alpha^2+\omega^4}.
\]

Elevando al cuadrado:

\[
\frac{a_0^2}{R^2}=\alpha^2+\omega^4.
\]

Por tanto:

\[
\boxed{
\alpha^2+\omega^4=C
},
\]

donde:

\[
C=\frac{a_0^2}{R^2}.
\]

Esta condición no obliga a que:

\[
\alpha=0.
\]

Por ejemplo, durante un intervalo en que \(\omega\) aumente
desde un valor positivo menor que \(C^{1/4}\), podemos elegir:

\[
\alpha=\sqrt{C-\omega^4}.
\]

Como:

\[
\alpha=\deriv{\omega}{t},
\]

esta elección describe una evolución de \(\omega(t)\) con
aceleración angular positiva y no nula.

La rapidez:

\[
v=R|\omega|
\]

aumenta con el tiempo, aunque:

\[
a=a_0
\]

permanece constante en magnitud.

Por tanto, no se trata de MCU.

El error conceptual de la afirmación inicial consiste en
confundir la magnitud de la aceleración con la condición
de rapidez constante.

El movimiento circular uniforme exige:

\[
v=\text{constante},
\]

equivalentemente, para radio fijo:

\[
|\omega|=\text{constante}.
\]

Una aceleración de magnitud constante, por sí sola, no
garantiza esta condición.
\end{explicacion}